\documentclass[10pt,a4paper]{amsart}
\usepackage[margin=19mm]{geometry}
\usepackage{amsmath,amssymb,mathtools}
\usepackage[scr=rsfs,cal=euler]{mathalfa}
\usepackage{xfrac}
\usepackage[unicode,hidelinks,bookmarks=false]{hyperref}
\newcommand{\ie}{i.e.\ }
\newcommand{\cf}{cf.\ }
\newcommand{\resp}{respectively\ }
\newcommand{\Subsection}{\subsection}
\providecommand{\nobreakdash}{\nobreak\hbox{-}}
\setlength{\emergencystretch}{2em}
\multlinegap0pt
\usepackage{bm}
\usepackage{bbm}
\usepackage{dsfont}
\usepackage{xspace}
\usepackage{cancel}
\usepackage{mathtools}
\usepackage{amsthm}
\newtheorem{theorem}{Theorem}
\newtheorem{proposition}[theorem]{Proposition}
\newcommand{\Om}{\Omega}
\title[Spatially Limited Immune Access Creates Tumour Refugia in He]{Spatially Limited Immune Access Creates Tumour Refugia in Hepatocellular Carcinoma}
\author{Jiguang Yu, Louis Shuo Wang}
\date{Source: arXiv:2609.27949v1 (CC BY 4.0)}
\begin{document}
\maketitle

\noindent\emph{Source and licence.} Spatially Limited Immune Access Creates Tumour Refugia in Hepatocellular Carcinoma. Version arXiv:2609.27949v1 (CC BY 4.0), distributed under the Creative Commons Attribution 4.0 International licence, as recorded in the arXiv licence metadata for this exact version and shown on the abstract page https://arxiv.org/abs/2609.27949v1. Full rights notice: licenses/arxiv-2609.27949-CC-BY-4.0.txt.\par\smallskip

\noindent\textit{Editorial scope.} Counted unit: the Proposition whose source label is \texttt{app:prop:positivity} in the article (source0.tex lines 642--645, source label \texttt{app:prop:positivity}), delivered together with the article's own complete original proof of it (source0.tex lines 647--675, one \texttt{proof} environment of 29 lines). No mathematics is written, repaired or re-derived: statement and proof are the article's own text line for line. Required context retained uncounted: none. Every definition, display and label of those lines is retained unchanged. Omitted from this package and not redistributed: 1--641, 646--646, 676--1026. Nothing inside a retained excerpt is omitted or altered: every retained body is delivered exactly as the article's lines are written, so no omission receipt of any kind accompanies this package. Citations: the retained excerpts carry no citation command at all, so no bibliography entry of the article is transcribed and no reference is redistributed. Reference adaptation: none was needed and none was made; every reference inside the delivered unit resolves inside it, so no unresolved reference or citation is produced. Printed slips kept literal: no printed word, symbol, hypothesis or number of the counted statement or of its proof was altered. Layout and class: the article's own class is \texttt{interact} with packages including epstopdf, subfig, float, natbib, xcolor; the wrapper is the lane's pinned \texttt{amsart} editorial wrapper at 10pt on A4, which restates the theorem-like environments the retained text needs and adds the article's own macro definitions that the retained text needs (\texttt{\textbackslash Om}), with extra wrapper packages (bm, bbm, dsfont, xspace, cancel, mathtools). The delivered numbering is the unit's own. Assets: no retained span contains a figure environment, a graphics-inclusion command, a PDF-page inclusion, a table or a TikZ picture; no image file is copied into this package, the package declares no asset dependency and no image file is redistributed with it. Licence: version arXiv:2609.27949v1 of this article is distributed under the Creative Commons Attribution 4.0 International licence, as recorded in the arXiv licence metadata for this exact version and shown on the abstract page https://arxiv.org/abs/2609.27949v1. A full rights notice is delivered at licenses/arxiv-2609.27949-CC-BY-4.0.txt. Characters: every retained excerpt is pure ASCII, so the delivered engine, XeLaTeX with its default Latin Modern text font, reports no missing character.
\begin{proposition}[Nonnegativity]\label{app:prop:positivity}
For nonnegative initial data the classical solution satisfies $u,v,w\ge0$ on
$[0,T_{\max})$.
\end{proposition}

\begin{proof}
Work with the truncated system in which $w/(1+w)$ and the $\gamma uv$ source use
$w^+$ and $v^+$; its right-hand side is locally Lipschitz, so the local theory
applies and it suffices to prove nonnegativity there. Since the $u$-reaction
$u(1-u-v)$ vanishes at $u=0$, the scalar comparison principle gives $u\ge0$. For
$v$, fix $T<T_{\max}$ and set $v_-=\max\{-v,0\}$. The map $z\mapsto\max\{-z,0\}$
is Lipschitz, so $v_-(\cdot,t)\in H^1(\Om)$ with
$\nabla v_-=-\mathbf 1_{\{v<0\}}\nabla v$ and
$t\mapsto\frac12\int_\Om v_-^2$ is absolutely continuous; all identities below
are in the weak sense. Because the solution is classical
$C^{2,1}$ on $\overline\Om\times[0,T]$, $\|\Delta w(\cdot,t)\|_{L^\infty}$ is
finite and bounded on $[0,T]$, which justifies the Gronwall step. Testing the
$v$-equation with $v_-$ and using the Neumann conditions,
\[
-\tfrac12\tfrac{d}{dt}\!\int_\Om v_-^2
= d_2\!\int_\Om|\nabla v_-|^2 +\tfrac{\xi}{2}\!\int_\Om(\Delta w)v_-^2
+\int_{\{v<0\}}\!\Big(\sigma_0+\sigma_1\tfrac{w^+}{1+w^+}-\delta v-\beta uv\Big)v_-,
\]
where the diffusion term $d_2\!\int_\Om|\nabla v_-|^2\ge0$, the truncated source
is $\ge0$, and on $\{v<0\}$ the terms $-\delta v,-\beta uv$ contribute
nonnegative multiples of $v_-$. Dropping these nonnegative contributions and
bounding the chemotaxis term by
$\tfrac{\xi}{2}\|\Delta w\|_{L^\infty}\!\int_\Om v_-^2$ gives
$\frac{d}{dt}\int_\Om v_-^2\le\xi\|\Delta w\|_{L^\infty}\int_\Om v_-^2$, and
since $v_-(\cdot,0)=0$, Gronwall gives $v_-\equiv0$, i.e.\ $v\ge0$. Finally
$w_t-d_3\Delta w+\ell w=\alpha u+\gamma uv\ge0$ with $w_0\ge0$ gives $w\ge0$ by
comparison. Then $v^+=v$, $w^+=w$, so the truncated and original systems
coincide; uniqueness finishes the proof.
\end{proof}

\end{document}
